% !TeX root = ../navier-stokes-manuscript.tex
% ============================================================
% 2.1 Vortex Stretching
% ============================================================

The momentum equation has four pieces:
\[
  \underbrace{\partial_t u}_{\text{local acceleration}}
  +
  \underbrace{(u\cdot\nabla)u}_{\text{nonlinear transport}}
  =
  \underbrace{-\nabla p}_{\text{pressure}}
  +
  \underbrace{\nu\Delta u}_{\text{viscous diffusion}}.
\]
The central regularity problem lies in the competition between nonlinear transport and viscosity.
The hope is simple: viscosity smooths the fluid faster than nonlinear transport can sharpen it.
If that remains true at every scale, the velocity stays smooth.

\subsection{Vortex Stretching}\label{sub:ThePerfectlyStretchingVortex}

In three dimensions, nonlinear transport can sharpen the velocity field by stretching vorticity.
Take one material vortex segment whose vorticity runs along its axis.
If \(e\) is the axial direction and
\[
  S:=\frac12\bigl(\nabla u+(\nabla u)^{\mathsf T}\bigr),
  \qquad
  Se=\sigma(t)e,
  \qquad
  \sigma(t)>0,
\]
then the surrounding strain pulls the segment out at the rate
\[
  \frac{\dot L(t)}{L(t)}
  =
  e\cdot Se
  =
  \sigma(t).
\]
The segment lengthens, but its material volume cannot change.
Writing \(\mathcal V=A(t)L(t)\) and \(r_{\mathrm{eff}}(t):=\sqrt{A(t)/\pi}\), incompressibility gives
\[
  \frac{d}{dt}(A L)=0,
  \qquad
  \frac{\dot A}{A}=-\sigma(t),
  \qquad
  \frac{\dot r_{\mathrm{eff}}}{r_{\mathrm{eff}}}
  =
  -\frac{\sigma(t)}{2}.
\]
The same strain that lengthens the segment is driving its radius down.

It acts on the rotation at the same time.
For \(\omega:=\nabla\times u\),
\[
  \frac{D\omega}{Dt}
  =
  S\omega+\nu\Delta\omega,
  \qquad
  \frac{D}{Dt}:=\partial_t+u\cdot\nabla.
\]
When \(\omega=|\omega|e\), the stretching contribution is
\[
  S\omega=\sigma(t)\omega,
  \qquad
  \frac12\frac{D|\omega|^2}{Dt}
  =
  \underbrace{\sigma(t)|\omega|^2}_{\text{stretching}}
  +
  \underbrace{\nu\,\omega\cdot\Delta\omega}_{\text{viscosity}}.
\]
Stretching pushes the rotation upward while viscosity opposes that growth.
In the perfectly stretching case, the positive stretching contribution persists instead of being cancelled by viscosity, so the rotation rises as the radius falls.

Finite kinetic energy does not stop this motion.
If \(|u|\leq U_0\) on the fixed-volume segment \(\Omega_{\mathrm{seg}}(t)\), then
\[
  E_{\mathrm{seg}}(t)
  :=
  \frac12\int_{\Omega_{\mathrm{seg}}(t)}|u(x,t)|^2\,dx
  \leq
  \frac12U_0^2\mathcal V
  <
  \infty,
  \qquad
  \frac{|\omega|}{\sqrt2}
  =
  \left|\frac12\bigl(\nabla u-(\nabla u)^{\mathsf T}\bigr)\right|_{\mathrm F}
  \leq
  |\nabla u|_{\mathrm F}.
\]
The segment can carry finite energy while it becomes longer, thinner, and faster rotating.
The gradients rise.

\divider{\NavierStokes{} Scaling}

The shrinking radius now puts the original hope to a scale test.
At a fixed time \(t_0\), let \(\ell:=r_{\mathrm{eff}}(t_0)\), and for \(\lambda>1\) define the rescaled solution
\[
  u_\lambda(x,t)
  :=
  \lambda u(\lambda x,\lambda^2t),
  \qquad
  p_\lambda(x,t)
  :=
  \lambda^2p(\lambda x,\lambda^2t).
\]
At time \(\lambda^{-2}t_0\), the corresponding structure in \(u_\lambda\) has width \(\lambda^{-1}\ell\).
The four terms at that finer width transform together:
\[
\begin{aligned}
  \partial_tu_\lambda(x,t)
  &=
  \lambda^3(\partial_tu)(\lambda x,\lambda^2t),
  \\
  (u_\lambda\cdot\nabla)u_\lambda(x,t)
  &=
  \lambda^3\bigl((u\cdot\nabla)u\bigr)(\lambda x,\lambda^2t),
  \\
  \nabla p_\lambda(x,t)
  &=
  \lambda^3(\nabla p)(\lambda x,\lambda^2t),
  \\
  \nu\Delta u_\lambda(x,t)
  &=
  \lambda^3\nu(\Delta u)(\lambda x,\lambda^2t).
\end{aligned}
\]
Moving to a narrower scale makes viscosity stronger, but it makes nonlinear sharpening stronger by exactly the same amount.
Viscosity gains no advantage from scale.

\divider{Critical Height}

That leaves the size of the velocity itself.
With \(\Lambda:=(-\Delta)^{1/2}\), the same rescaling gives
\[
  \norm{u_\lambda(t)}_{\dot H^s}^2
  =
  \lambda^{2s-1}
  \norm{u(\lambda^2t)}_{\dot H^s}^2.
\]
Below \(s=\tfrac12\) the norm falls as the structure concentrates; above it the norm rises.
At \(s=\tfrac12\), scale contributes no factor at all.
Set
\[
  H(t)
  :=
  \frac12\norm{u(t)}_{\dot H^{1/2}}^2
  =
  \frac12\norm{\Lambda^{1/2}u(t)}_2^2.
\]
The neighboring levels separate while this one remains fixed:
\[
  \norm{u_\lambda(t)}_2^2
  =
  \lambda^{-1}\norm{u(\lambda^2t)}_2^2,
  \qquad
  H_\lambda(t)=H(\lambda^2t),
  \qquad
  \norm{\nabla u_\lambda(t)}_2^2
  =
  \lambda\norm{\nabla u(\lambda^2t)}_2^2.
\]
Concentration can lower kinetic energy and raise first-derivative energy without moving \(H\).
Any rise or fall in the critical height must come from the equation, not from the change of scale.

\divider{Nonlinear Production versus Viscous Dissipation}

Pairing the momentum equation at this height exposes the same contest directly.
Write its signed nonlinear contribution as
\[
  \mathcal P_{1/2}(t)
  :=
  -\operatorname{Re}
  \pair
    {\Lambda^{1/2}u(t)}
    {\Lambda^{1/2}\bigl((u(t)\cdot\nabla)u(t)\bigr)}_{L^2}.
\]
Incompressibility removes the pressure pairing, and the Laplacian gives
\[
  H'(t)
  +
  \nu\norm{\Lambda^{3/2}u(t)}_2^2
  =
  \mathcal P_{1/2}(t).
\]
The critical height rises when nonlinear production exceeds viscous dissipation and falls when viscosity wins.
At a finer scale, both rates acquire the same factor:
\[
  \mathcal P_{1/2}[u_\lambda](t)
  =
  \lambda^2\mathcal P_{1/2}[u](\lambda^2t),
  \qquad
  \nu\norm{\Lambda^{3/2}u_\lambda(t)}_2^2
  =
  \lambda^2\nu\norm{\Lambda^{3/2}u(\lambda^2t)}_2^2.
\]
The balance still does not tilt toward viscosity.
Instead, the entire contest has moved onto the shorter time coordinate \(t\mapsto\lambda^{-2}t\).

\divider{The Heat Clock}

Viscosity carries that same time scale.
For the linear heat equation \(v_t=\nu\Delta v\),
\[
  \widehat v(\xi,t+\delta t)
  =
  e^{-\nu|\xi|^2\delta t}\widehat v(\xi,t).
\]
A structure of width \(r\) lives around frequencies \(|\xi|\simeq r^{-1}\).
Diffusion changes it by order one when \(\nu|\xi|^2\delta t\simeq1\), which fixes the viscous clock
\[
  \tau_\nu(r)
  :=
  \frac{r^2}{\nu}.
\]
Shrinking the width by \(\lambda^{-1}\) gives
\[
  \tau_\nu(\lambda^{-1}r)
  =
  \lambda^{-2}\tau_\nu(r).
\]
The same square that shortens the \NavierStokes{} time coordinate shortens the time viscosity has to act at the finer width.
Each smaller-scale comparison is faster, but the sharpening and the smoothing remain locked to the same scale.

\divider{The Zeno Cascade}

Let successive widths halve:
\[
  r_n:=2^{-n}r_0,
  \qquad
  \tau_n:=\frac{r_n^2}{\nu}.
\]
Their viscous clocks form a summable sequence:
\[
  \tau_n
  =
  \frac{r_0^2}{\nu}4^{-n},
  \qquad
  \sum_{n=0}^{\infty}\tau_n
  =
  \frac{4r_0^2}{3\nu}
  <
  \infty.
\]
Infinitely many scale comparisons fit inside a finite total time.
That sum does not settle the competition; it shows how little time remains in which the same competition must be settled.

A singular cascade would have to occur in one solution, through widths
\[
  r_0>r_1>r_2>\cdots\longrightarrow0
\]
at times
\[
  t_0<t_1<t_2<\cdots\uparrow T_*<\infty,
  \qquad
  t_{n+1}-t_n\asymp\tau_n,
\]
with comparison constants independent of \(n\).
The remaining tail would then satisfy
\[
  T_*-t_n
  \asymp
  \sum_{k=n}^{\infty}\tau_k
  =
  \frac{4r_n^2}{3\nu}.
\]
At every finer width, the same velocity field has only an interval of the order of that width's viscous clock in which to change again.
Spatial concentration has become a demand on temporal response.

Assume this same solution loses continuation at \(T_*\) through blow-up in one fixed continuation-controlling Banach space \(B\).
If some finite temporal row \(\partial_t^N u\), \(N\geq1\), remained uniformly bounded in \(B\) after a time \(t_0<T_*\), then
\[
  \partial_t^{N-1}u(t)
  =
  \partial_t^{N-1}u(t_0)
  +
  \int_{t_0}^{t}\partial_t^N u(s)\,ds
\]
would remain bounded on the finite interval.
Repeating the same integration down through the lower rows would keep \(u\) bounded in \(B\), contradicting the assumed blow-up.
No finite temporal row can retain uniform control along such a singular history.

The first row is the velocity, the momentum equation supplies its acceleration, and each further time differentiation supplies the next response:
\[
  u,
  \qquad
  \partial_tu,
  \qquad
  \partial_t^2u,
  \qquad
  \partial_t^3u,
  \qquad
  \ldots
\]
Together they form the temporal Tower.
